% ============================================================
% Add this to your preamble
% ============================================================
\definecolor{promptbg}{gray}{0.95}
\definecolor{promptframe}{gray}{0.5}
\definecolor{afterbg}{HTML}{F0FAF0}
\definecolor{afterframe}{HTML}{2E7D32}
\definecolor{beforebg}{HTML}{FFF5F5}
\definecolor{beforeframe}{HTML}{C62828}
\definecolor{thinkbg}{HTML}{FAFAFA}
\definecolor{thinkframe}{gray}{0.75}

% --- mdframed style definitions ---
\mdfdefinestyle{promptbox}{%
  backgroundcolor=promptbg,
  linecolor=promptframe,
  linewidth=0.8pt,
  innertopmargin=8pt,
  innerbottommargin=8pt,
  innerleftmargin=8pt,
  innerrightmargin=8pt,
  skipabove=6pt,
  skipbelow=4pt
}
\mdfdefinestyle{afterbox}{%
  backgroundcolor=afterbg,
  linecolor=afterframe,
  linewidth=1pt,
  innertopmargin=8pt,
  innerbottommargin=8pt,
  innerleftmargin=8pt,
  innerrightmargin=8pt,
  skipabove=6pt,
  skipbelow=4pt
}
\mdfdefinestyle{beforebox}{%
  backgroundcolor=beforebg,
  linecolor=beforeframe,
  linewidth=1pt,
  innertopmargin=8pt,
  innerbottommargin=8pt,
  innerleftmargin=8pt,
  innerrightmargin=8pt,
  skipabove=6pt,
  skipbelow=4pt
}
\mdfdefinestyle{thinkbox}{%
  backgroundcolor=thinkbg,
  linecolor=thinkframe,
  linewidth=0.5pt,
  innertopmargin=6pt,
  innerbottommargin=6pt,
  innerleftmargin=6pt,
  innerrightmargin=6pt,
  skipabove=4pt,
  skipbelow=4pt
}


% ============================================================
% Place the following in the body of the document
% ============================================================

\begin{figure*}[htbp]
\scriptsize
% === Full-width PROMPT ===
\begin{mdframed}[style=promptbox]
\noindent\textbf{프롬프트}

\vspace{2pt}\noindent
\textbf{System:} 다음 Sudoku 퍼즐을 푸세요. 빈 칸은 0으로 표시됩니다. 정답으로는 풀린 퍼즐만 공백으로 구분된 숫자의 9x9 그리드 형태로 즉시 출력하세요. \#\#\#\#, 설명, 그 밖의 어떤 텍스트도 포함하지 마세요.

\vspace{4pt}\noindent
\textbf{User:} 4 0 0 0 0 1 0 0 6 0 2 1 8 0 9 0 0 0 0 9 0 4 0 2 0 0 0 0 0 0 0 4 8 6 0 0 0 0 5 0 2 7 4 1 9 9 0 7 5 1 0 3 0 0 6 1 3 2 9 0 5 0 0 5 0 9 0 8 0 0 6 4 2 8 4 6 7 0 1 9 0
\end{mdframed}
\vspace{2pt}
% === Two columns: Before \sdrlfootnotesize | After \sdrlfootnotesize ===
\noindent
\begin{minipage}[t]{0.49\textwidth}
% --- Before \sdrlfootnotesize ---
\begin{mdframed}[style=beforebox]
\noindent\textbf{Downstream Finetuning 이전} \hfill {\colorbox{beforebg}{\textcolor{beforeframe}{\textbf{$\boldsymbol\times$~오답 }}}}

\noindent\textit{사고 과정 (일부 생략):}
\begin{mdframed}[style=thinkbox]
\texttt{<|turn>model} \\
\texttt{<|channel>thought} \\
\texttt{<channel|>}4 8 3 7 5 1 9 2 6 \\
7 2 1 8 6 9 5 3 4 \\
6 9 5 4 3 2 8 1 7 \\
1 4 8 9 9 6 3 7 5 \\
3 5 6 2 7 4 1 9 8 \\
9 7 2 5 1 3 4 6 1 \\
2 3 9 1 4 5 7 8 6 \\
5 6 7 1 2 8 4 4 3 \\
8 1 4 6 7 6 2 9 5 \\
\\
잠깐, 입력 문자열이 약간 조각나 있었으니 더 신중하게 다시 풀어 보겠습니다. \\
입력: \\
4 0 0 | 0 0 1 | 0 0 6 \\
0 2 1 | 8 0 9 | 0 0 0 \\
0 9 0 | 4 0 2 | 0 0 0 \\
0 0 0 | 0

{\itshape [\ldots 일부 생략]}
\end{mdframed}
\vspace{2pt}
\noindent\textit{응답:}
\begin{mdframed}[style=thinkbox]
\textit{[비어 있음 --- 답변이 잘림]}
\end{mdframed}
\end{mdframed}
\end{minipage}%
\hfill
\begin{minipage}[t]{0.49\textwidth}
% --- After \sdrlfootnotesize ---
\begin{mdframed}[style=afterbox]
\noindent\textbf{Downstream Finetuning 이후} \hfill {\colorbox{afterbg}{\textcolor{afterframe}{\textbf{\checkmark~정답}}}}

\noindent\textit{사고 과정:}
\begin{mdframed}[style=thinkbox]
\textit{[비어 있음 --- 사고 채널을 성공적으로 억제함]}
\end{mdframed}
\vspace{2pt}
\noindent\textit{응답:}
\begin{mdframed}[style=thinkbox]
4 5 8 7 3 1 9 2 6 \\
3 2 1 8 6 9 7 4 5 \\
7 9 6 4 5 2 8 3 1 \\
1 3 2 9 4 8 6 5 7 \\
8 6 5 3 2 7 4 1 9 \\
9 4 7 5 1 6 3 8 2 \\
6 1 3 2 9 4 5 7 8 \\
5 7 9 1 8 3 2 6 4 \\
2 8 4 6 7 5 1 9 3
\end{mdframed}
\end{mdframed}
\end{minipage}
\caption{\textbf{Open-Source Downstream Finetuning 전후의 Sudoku 퍼즐에서 서식 규칙 준수 여부를 평가한 예시.} Open-Source Downstream Finetuning 이전에는 모델이 엄격한 부정 제약(``풀린 퍼즐만 즉시 출력하라'')을 따르지 못하고, 사고 채널에서 잘못된 그리드를 구성하려 시도한 뒤 최종 응답을 제공하지 못한다. Finetuning 이후에는 모델이 응답 채널에 올바르게 풀린 9x9 그리드를 직접 출력한다.}
\label{fig:sudoku-rl-example}
\end{figure*}
